\documentclass[UTF8]{ctexart}
\usepackage[T1]{fontenc}
\usepackage{amsmath,amssymb}%數學公式、符號
\usepackage[top=2.5truecm,bottom=2.5truecm,
left=3truecm,right=2.5truecm]{geometry}%版面設置
\usepackage{graphicx, subfig}%圖形
\usepackage[colorlinks,linkcolor=purple,
anchorcolor=blue,citecolor=green]{hyperref}%超連結、pdf書籤
\usepackage{caption}
\usepackage{graphicx, subfig}%圖片宏包
\usepackage{type1cm} %調整字體絕對大小
\usepackage{textpos} %設定文字絕對位置
%\usepackage{fontspec}%加這個即可設定字體
\usepackage{xeCJK} %使用xeCJK，其他的還有CJK或是xCJK


%%%%%%%%%%%%%自訂字體%%%%%%%%%%%%%

\newcommand{\thirty}{\fontsize{30pt}\baselineskip\selectfont}%字體大小30pt
\newcommand{\twentyfour}{\fontsize{24pt}{\baselineskip}\selectfont}%字體大小24pt
\newcommand{\twenty}{\fontsize{20pt}{\baselineskip}\selectfont}%字體大小20pt
\newcommand{\eigh}{\fontsize{18pt}{\baselineskip}\selectfont}%字體大小18pt
\newcommand{\sixteen}{\fontsize{16pt}{\baselineskip}\selectfont}%字體大小16pt
\newcommand{\fourteen}{\fontsize{14pt}{\baselineskip}\selectfont}%字體大小14pt
\newcommand{\twelve}{\fontsize{12pt}{\baselineskip}\selectfont}%字體大小12pt

%%%%%%%%%%%%%正文開始%%%%%%%%%%%%%

\begin{document}
    \pagestyle{plain} %設定每頁頁碼
    {\renewcommand\baselinestretch{1.8}\selectfont %設定以下行距
    {\begin{center} %以下文字置中
        \kaishu \twentyfour{國立虎尾科技大學}\\{機械設計工程系}\\{專題製作報告}\\
        \hspace*{\fill} \\
        \bf \kaishu \thirty {機構設計與模擬實作}\\
    \end{center}}
    \par} %結束指定行距

     {\renewcommand\baselinestretch{1.2}\selectfont %設定以下行距
    {\begin{textblock}{30}(2,2.5) %{寬度}(以左上角為原點之右移量,下移量)
    \noindent \kaishu \eigh \makebox[6em][s]{指導教授}\enspace:\qquad
    \kaishu \eigh \makebox[6em][s]{李武鉦}\\ %\noindent指定首行不進行縮排
    \kaishu \eigh \makebox[6em][s]{班級}\enspace:\qquad
    \kaishu \eigh \makebox[6em][s]{四設計三甲} \\ %\makebox為文本盒子
    \kaishu \eigh \makebox[6em][s]{學生}\enspace:\qquad
    {\begin{minipage}[t][9.8em][s]{12em}
    \kaishu \eigh \setlength{\baselineskip}{10pt plus 2pt}
    林昱秀\qquad40723102\vfill 林晏瑩\qquad40723103\vfill 劉光智\qquad40723145\vfill 吳佳穎\qquad40723153\vfill 蔡育灃\qquad40723245
    \end{minipage}}
    \end{textblock}} %指定以下文字整體偏移至頁面的絕對位置
    \par} %結束指定行距

    {\begin{textblock}{12}(0,8)
    {\begin{center}
    \noindent \kaishu \eigh \makebox[12em][s]{中華民國一一零年一月}
    \end{center}}
    \end{textblock}}
    \newpage
%%%%%%%%%%%%%%%%page2%%%%%%%%%%%%%%%%
    {\begin{center}
    {\renewcommand \baselinestretch {1.4}  \kaishu \eigh {國立虎尾科技大學 \qquad 機械設計工程系}\\
    {學生專題製作合格認可證明}\\
    \hspace*{\fill} \\ %似enter鍵換行
    \par}
     \end{center}}
    {\begin{textblock}{60}(0.7,0)
    \noindent \kaishu \fourteen 專題製作修習學生\enspace:\quad
    {\begin{minipage}[t]{10em}\underline{　　　　　　　　　　　　}\\ \underline{　　　　　　　　　　　　}\\ \underline{　　　　　　　　　　　　}\\ \underline{　　　　　　　　　　　　}\\ \underline{　　　　　　　　　　　　}\\ %下劃線符號指令
    \end{minipage}}
    {\begin{textblock}{30}(0,0.3)
    \noindent \kaishu \fourteen 專題製作題目\enspace:\quad
    \hspace*{\fill} \\
    \hspace*{\fill} \\
    \noindent \kaishu \fourteen 經評量合格，特此證明
    \hspace*{\fill} \\
    \noindent \kaishu \fourteen \makebox[6em][s]{評審委員}\enspace:\quad
    {\begin{minipage}[t]{6em} \underline{　　　　　　　　　　　　}\\ \underline{　　　　　　　　　　　　}\\ \underline{　　　　　　　　　　　　}\\
    \end{minipage}}
    \end{textblock}}
    {\begin{textblock}{10}(0,6.5)
    {\begin{flushleft}
    \kaishu \fourteen \makebox[6em][s]{指導老師}\enspace:\quad \underline{　　　　　　　　　　　　}\\
    \kaishu \fourteen \makebox[6em][s]{系主任}\enspace:\quad \underline{　　　　　　　　　　　　}\\
    \hspace*{\fill} \\
    \kaishu \fourteen \makebox[6em][s]{中華民國}\qquad\quad{年}\qquad\quad{月}\qquad\quad{日}
    \end{flushleft}}
    \end{textblock}}
    \end{textblock}}
    \newpage
%%%%%%%%%%%%%%%%page3%%%%%%%%%%%%%%%%
    {\begin{center}
        \bf \kaishu \twentyfour {摘要}\\
    \end{center}}
        \kaishu \sixteen 生產自動化是現今工業界中最重要的一環，一國國力之強弱取決於工業發展的情況，而其經濟的動脈更是來自於自動化工業，故發展自動化工業技術，已是時代潮流。對於我國在踏入高科技產業的同時，如何以更低的成本與縮短生產製程，來提高在國際上的競爭力，是大家所努力的目標，而身為未來二十一世紀的一員，更需了解其重要性。
        
        目前於市面上的數位控制加工機（CNC），其成本昂貴，且體積龐大，並不方便於一般教學的使用，導致技術無法有效的傳授與教導，為了使教學上能更便利於授課，且學生能立即操作而了解其動作原理，故本組決定運用在校所學之相關課程，以完成一部具有高精度、體積小與低成本（十五萬元以下）的 PC-Based 三軸運動控制實驗台為研究目標，達成在教學上的需要，並增加系內的實驗設備。
        
        本專題主要目的在了解伺服馬達的控制原理、三軸運動控制卡的使用方式、人機介面程式的撰寫與增加機械實務加工的能力。
	\newpage
%%%%%%%%%%%%%%%%page4%%%%%%%%%%%%%%%%
	{\begin{center}
        \bf \kaishu \twentyfour {作者簡介}\\
    \end{center}}
    {\begin{textblock}{6}(0,0.5)
    \begin{figure}
        \includegraphics[width=1.15in]{st1.jpg} %{}內是圖片文件的相對路徑
    \end{figure}
    \end{textblock}}
    {\renewcommand\baselinestretch{0.99}\selectfont %設定以下行距
    {\begin{textblock}{15}(3.5,1) %{寬度}(以左上角為原點之右移量,下移量)
    \noindent \kaishu \fourteen \makebox[4em][s]{姓名}\enspace:\enspace
    \kaishu \fourteen \makebox[4em][s]{林昱秀}\\ %\noindent指定首行不進行縮排
    \kaishu \fourteen \makebox[4em][s]{學號}\enspace:\enspace
    \kaishu \fourteen \makebox[4em][s]{40723102} \\ %\makebox為文本盒子
    \kaishu \fourteen \makebox[4em][s]{畢業學校}\enspace:\enspace
    \kaishu \fourteen \makebox[9em][s]{國立虎尾科技大學}\\
    \kaishu \fourteen \makebox[5em][s]{\quad}\enspace\enspace
    \kaishu \fourteen \makebox[8em][s]{機械設計工程系}\\
    \kaishu \fourteen \makebox[4em][s]{經歷}\enspace:\enspace
    \end{textblock}}}
    \hspace*{\fill} \\
    
    {\begin{textblock}{6}(0,2.15)
    \begin{figure}
        \includegraphics[width=1.15in]{st2.jpg} %{}內是圖片文件的相對路徑
    \end{figure}
    \end{textblock}}
    {\renewcommand\baselinestretch{0.99}\selectfont %設定以下行距
    {\begin{textblock}{15}(3.5,2.65) %{寬度}(以左上角為原點之右移量,下移量)
    \noindent \kaishu \fourteen \makebox[4em][s]{姓名}\enspace:\enspace
    \kaishu \fourteen \makebox[4em][s]{林晏瑩}\\ %\noindent指定首行不進行縮排
    \kaishu \fourteen \makebox[4em][s]{學號}\enspace:\enspace
    \kaishu \fourteen \makebox[4em][s]{40723103} \\ %\makebox為文本盒子
    \kaishu \fourteen \makebox[4em][s]{畢業學校}\enspace:\enspace
    \kaishu \fourteen \makebox[9em][s]{國立虎尾科技大學}\\
    \kaishu \fourteen \makebox[5em][s]{\quad}\enspace\enspace
    \kaishu \fourteen \makebox[8em][s]{機械設計工程系}\\
    \kaishu \fourteen \makebox[4em][s]{經歷}\enspace:\enspace
    \end{textblock}}}
    \hspace*{\fill} \\
    

    {\begin{textblock}{6}(0,3.8)
    \begin{figure}
        \includegraphics[width=1.15in]{st3.jpg} %{}內是圖片文件的相對路徑
    \end{figure}
    \end{textblock}}
    {\renewcommand\baselinestretch{0.99}\selectfont %設定以下行距
    {\begin{textblock}{15}(3.5,4.3) %{寬度}(以左上角為原點之右移量,下移量)
    \noindent \kaishu \fourteen \makebox[4em][s]{姓名}\enspace:\enspace
    \kaishu \fourteen \makebox[4em][s]{劉光智}\\ %\noindent指定首行不進行縮排
    \kaishu \fourteen \makebox[4em][s]{學號}\enspace:\enspace
    \kaishu \fourteen \makebox[4em][s]{40723145} \\ %\makebox為文本盒子
    \kaishu \fourteen \makebox[4em][s]{畢業學校}\enspace:\enspace
    \kaishu \fourteen \makebox[9em][s]{國立虎尾科技大學}\\
    \kaishu \fourteen \makebox[5em][s]{\quad}\enspace\enspace
    \kaishu \fourteen \makebox[8em][s]{機械設計工程系}\\
    \kaishu \fourteen \makebox[4em][s]{經歷}\enspace:\enspace
    \end{textblock}}}
    \hspace*{\fill} \\
    
    {\begin{textblock}{6}(0,5.45)
    \begin{figure}
        \includegraphics[width=1.15in]{st4.jpg} %{}內是圖片文件的相對路徑
    \end{figure}
    \end{textblock}}
    {\renewcommand\baselinestretch{0.99}\selectfont %設定以下行距
    {\begin{textblock}{15}(3.5,5.95) %{寬度}(以左上角為原點之右移量,下移量)
    \noindent \kaishu \fourteen \makebox[4em][s]{姓名}\enspace:\enspace
    \kaishu \fourteen \makebox[4em][s]{吳佳穎}\\ %\noindent指定首行不進行縮排
    \kaishu \fourteen \makebox[4em][s]{學號}\enspace:\enspace
    \kaishu \fourteen \makebox[4em][s]{40723153} \\ %\makebox為文本盒子
    \kaishu \fourteen \makebox[4em][s]{畢業學校}\enspace:\enspace
    \kaishu \fourteen \makebox[9em][s]{國立虎尾科技大學}\\
    \kaishu \fourteen \makebox[5em][s]{\quad}\enspace\enspace
    \kaishu \fourteen \makebox[8em][s]{機械設計工程系}\\
    \kaishu \fourteen \makebox[4em][s]{經歷}\enspace:\enspace
    \end{textblock}}}
    \hspace*{\fill} \\
    
    {\begin{textblock}{6}(0,7.1)
    \begin{figure}
        \includegraphics[width=1.15in]{st5.jpg} %{}內是圖片文件的相對路徑
    \end{figure}
    \end{textblock}}
    \renewcommand\baselinestretch{0.99}\selectfont %設定以下行距
    {\begin{textblock}{15}(3.5,7.6) %{寬度}(以左上角為原點之右移量,下移量)
    \noindent \kaishu \fourteen \makebox[4em][s]{姓名}\enspace:\enspace
    \kaishu \fourteen \makebox[4em][s]{蔡育灃}\\ %\noindent指定首行不進行縮排
    \kaishu \fourteen \makebox[4em][s]{學號}\enspace:\enspace
    \kaishu \fourteen \makebox[4em][s]{40723245} \\ %\makebox為文本盒子
    \kaishu \fourteen \makebox[4em][s]{畢業學校}\enspace:\enspace
    \kaishu \fourteen \makebox[9em][s]{國立虎尾科技大學}\\
    \kaishu \fourteen \makebox[5em][s]{\quad}\enspace\enspace
    \kaishu \fourteen \makebox[8em][s]{機械設計工程系}\\
    \kaishu \fourteen \makebox[4em][s]{經歷}\enspace:\enspace
    \end{textblock}}
    \par
    \hspace*{\fill} \\
    \newpage
%%%%%%%%%%%%%%%%page5%%%%%%%%%%%%%%%%
    
\end{document} 